\chapter{How the product works}
\label{ch:design}

When an author asks for a humanized revision, the team should be able to follow
the work from the immutable source to the final reader decision. Humanvoice keeps
that correspondence between the intended audience, every substantive source
concept, each rewritten passage, the verification checks, the source and its
rendering, and the outcome. Once the correspondence is stable, engineers can
replace a parser or model without losing the story of why a paragraph says what
it says or why it was released.

\section{The flow of one revision}

The sequence follows the work: state the reader's task, inventory and freeze the
source concepts, rewrite passages in dependency order, verify the candidate,
repair located failures, and ask the reader only after release.

\begin{figure}[H]
\centering
\resizebox{\textwidth}{!}{%
\begin{tikzpicture}[node distance=6mm and 4mm,font=\small]
  \node[hvstage,text width=25mm] (brief) {Humanization brief\\reader, genre, voice};
  \node[hvstage,right=of brief,text width=25mm] (plan) {Frozen baseline\\concepts, obligations, dependencies};
  \node[hvgoldbox,right=of plan,text width=25mm] (draft) {Source rewrite\\one passage at a time};
  \node[hvwarn,right=of draft,text width=25mm] (preflight) {Semantic preflight\\retention and explanation};
  \node[hvbluebox,right=of preflight,text width=25mm] (release) {Release gate\\clean packet or stop};
  \node[hvbluebox,right=of release,text width=25mm] (read) {Human reader\\decision and action};
  \draw[hvflow] (brief)--(plan)--(draft)--(preflight)--(release)--(read);
  \draw[hvback] (preflight.south) to[bend right=25] node[below,hvnote]{located repair} (draft.south);
  \node[below=9mm of plan,align=center,text=hvslate,text width=34mm]
    (prov) {provenance: brief, baseline, source and render hashes};
  \node[below=9mm of release,align=center,text=hvslate,text width=34mm]
    (objects) {protected objects: equations, numbers, citations, claims};
  \draw[hvflow,dashed] (plan.south)--(prov.north);
  \draw[hvflow,dashed] (release.south)--(objects.north);
  \node[hvguard,fit=(preflight)(release),inner sep=3mm] {};
\end{tikzpicture}
}%
\caption{The common record keeps a defective revision inside a bounded repair
loop. Only a passed release gate reaches the expensive human reader.}
\label{fig:product-architecture}
\end{figure}

The brief comes first because a finding means nothing without a purpose. The
frozen baseline comes before any rewriting because a fluent paragraph can still
drop a distinction, mention a concept without explaining it, or use an idea
before its prerequisite. The source snapshot and protected spans are attached to
each unit so that a repair remains reproducible. The reader decision comes last
because it belongs to the released document, not to a model score or an internal
queue.

\section{The document representation}

The system has one source format: LaTeX. The common representation can later
admit Markdown or a word-processing adapter, but that extension is earned only
after locations and protected objects are preserved in LaTeX. The record model
starts before the source file. Each object has one clear owner and one reason it
exists in the path.

\begin{table}[H]
\centering
\small
\begin{tabularx}{\textwidth}{@{}p{0.20\textwidth}L p{0.25\textwidth}@{}}
\toprule
Object & Required contents & Owner of the decision \\
\midrule
HumanizationBrief & Intended reader, genre, known vocabulary, named exemplars,
evidence boundary, protected-object policy, voice constraints, privacy class,
remote-inference authorization, and unresolved assumptions. & Owner confirms that
a missing answer is safe to carry forward. \\
ConceptBaseline & Complete span partition, substantive concepts and their
teaching roles, explanation obligations, dependency order, scaffolding
dispositions, bidirectional coverage reviews, and the named freeze decision. &
Author and domain reviewer adjudicate ambiguity and sign the baseline. \\
DocumentVersion & Source hash, source format, rendered hash, author,
humanization brief, creation time, and disclosure status. & Technical owner
confirms the snapshot. \\
SemanticPreflightResult & Rule and critic versions, build result, span coverage,
concept correspondence, obligation fulfillment, dependency order, unsupported
additions, scaffolding leakage, exact-object integrity, and abstentions. &
Preflight controller records pass, repair, or fail-closed. \\
ProtectedSpan & Stable identifier, source range, object type, normalized
content hash, and allowed transformations. & Author declares what must be
preserved. \\
Finding & Source range, category, evidence, context, confidence or rule
version, editorial question, and disposition. & Author or editor accepts,
rejects, or edits. \\
RevisionRecord & Parent version, changed ranges, accepted findings, rejected
findings, and explanation of substantive changes. & Author owns the wording. \\
ReleaseDecision & Gate results, packet hashes, blocking reasons, and the person
who authorized release. & System blocks by default; there is no generic release
exception. \\
ReaderDecision & Reader role, document hash, rubric responses, elapsed time,
acceptance or requested changes, and confidence. & Named reader accepts or
returns the document. \\
\bottomrule
\end{tabularx}
\caption{The record model connects intent, source meaning, prose, checks, and
reader judgment. It is deliberately small enough to inspect and rich enough to
stop a lossy revision before release.}
\label{tab:record-model}
\end{table}

Protected objects include equations, numerical literals, units, labels,
citations, quotations, tables, code blocks, declared assumptions, and
load-bearing claims. The frozen baseline adds a second, stronger class of
protection: every substantive concept must survive with verified correspondence,
and prerequisites must be taught before dependents. Exact-object protection
does not make an object immutable. It makes a change visible and requires an
explanation. An author may intentionally change an estimate or replace a
citation; the system must distinguish that action from silent loss.

\section{The pre-human release contract}

The product is successful only if it changes when a human is asked to read. A
human may inspect a brief, resolve a high-consequence ambiguity, or adjudicate a
meaning question during construction. The expensive reader should not be asked
to find ordinary defects that the system could have found before release.

Every candidate version passes four gates:

\begin{enumerate}
\item \textbf{Integrity.} The source builds reproducibly; all protected spans
      are located or explicitly marked unparsed; citations and references
      resolve; and no seeded change to a number, unit, sign, equation, table,
      quotation, or claim scope can pass silently.
\item \textbf{Meaning.} Every frozen source concept has verified correspondence
      in the revision; every active explanation obligation is discharged with
      located evidence; every load-bearing claim keeps its evidence anchor; no
      unsupported claim is added; and every concept is taught before it is used
      or carries an explicit exemption.
\item \textbf{Writing.} Rewritten passages put concrete actors and objects before
      abstractions, keep the declared register, and vary sentence and paragraph
      pace. These are calibrated checks against named exemplars, not a universal
      style score, and never a length target.
\item \textbf{Reconstruction.} Isolated critics, calibrated on held-out human
      judgements, attempt to recover the purpose, mechanism, evidence,
      qualification, and requested action. They must report uncertainty and
      abstain when the evidence is insufficient.
\end{enumerate}

The release rule is fail-closed: a failed or missing gate produces a located
repair task, not a human packet. Repair has a bounded cycle count and an
oscillation detector. If the same defect returns, or two repairs undo one
another, the system stops and names the unresolved choice for the author. This
does not certify that the prose is beautiful. It prevents predictable failure
from being purchased with an executive's time.

A reviewer may correct and refreeze the baseline, or deactivate an obligation that
does not apply to this reader, with a recorded reason and lineage, before
verification runs. That is the only adjudication available, and it is not a
release exception. An unfrozen baseline, a source concept without verified
correspondence, an unmet active obligation, unresolved protected correspondence,
an unparsed promised object, unauthorized external transmission, missing
load-bearing evidence, and a broken build can never be excepted. The full
authority and record fields are fixed in
Appendix~\ref{app:implementation-contract}.

\begin{takeawaybox}{The human-attention budget}
The ordinary production path asks a human for a brief confirmation and one
final reading. It escalates earlier only for a decision the machine cannot own:
the intended meaning, a disputed source, or an explicit exception. The product
measures those escalations, rather than quietly treating them as free.
\end{takeawaybox}

\section{The diagnostic layer}

The first release needs only diagnostics that support a real author decision.

\begin{description}
\item[Register leakage.] Find private identifiers, file paths, status labels,
  and dates used as unexplained scope markers. Emit the source span and a
  question about the reader's missing context.
\item[Structural repetition.] Measure repeated section and paragraph shapes,
  opener distributions, and transition phrases against a genre reference.
  Report patterns for inspection; never impose a universal frequency target.
\item[Defensive or unsupported specificity.] Identify negative definitions,
  unsupported actor labels, and claims whose evidence or mechanism is not
  visible in the current document. These are candidate questions, not automatic
  edits.
\item[Paragraph packing and conceptual pacing.] Count the mechanisms,
  citations, and first-introduced concepts carried by each prose paragraph.
  Compare the target with the exemplars named in the reading brief and return
  ranked passages, not a universal score.
\item[Dependency order.] Check the first use of a term against known
  vocabulary, an introduction signal, a taught-later pointer, and a declared
  exemption zone. An uncertain match is an abstention for the author.
\item[Citation identity and build.] Resolve bibliography identity, unresolved
  references, and compilation failures. Queue sentence-to-source support for
  human verification.
\end{description}

During preflight, the comparison also looks for newly added numbers,
superlatives, comparisons, and causal verbs. An addition with no citation,
elementary derivation, or pointer to an existing protected claim becomes a
located assertion question. This catches unsupported specificity introduced
during an explanation pass without pretending that the detector can establish
the truth of the new sentence.

The critics are not one general-purpose judge. One checks concept order, one
checks evidence and scope, one checks concrete actors and mechanism, and one
attempts a low-context reconstruction. Their prompts, calibration set, order,
and disagreements are recorded. A critic that cannot beat a frozen baseline on
held-out examples is removed from the release path.

The system deliberately excludes per-document AI detection, population drift
monitoring, summarization, and a scalar acceptance score. Authorized
source-grounded rewriting and cause-specific repair are core product behavior;
autonomous acceptance is not. These exclusions keep the decision measurable
without asking a model to certify its own prose.

\section{The author and reader experience}

The author opens one session and sees the brief, the concept baseline, the source,
and the rendered revision together. Findings are grouped by location and gate,
with the most consequential blockers first. Selecting a finding shows its
evidence, the concept, obligation, or protected object it concerns, and the repair
it will trigger. The author can revise in the existing editor, or ask the writer
to rewrite the passage again; Humanvoice rebuilds and reverifies the result before
it advances.

The reader never sees the internal finding identifiers by default. The reader
receives the rendered document, the declared decision question, and a short
rubric only after the release gate passes. A reader may request a change without
diagnosing its cause. That is important: the product should not force the
funder to become a copy editor in order to provide useful feedback, nor should
it use that funder as a routine preflight service.

The session record keeps the two views linked. It can later answer questions
such as: which finding led to this sentence change; did the equation survive;
how long did the author spend resolving the queue; and did the reader accept
the resulting version? It cannot answer whether a different wording would have
been better unless the team runs that comparison.

\section{Privacy and disclosure}

Technical documents may contain unpublished estimates, client information, or
internal decisions. The default path is local parsing, local comparison, and
local storage. A remote model call is optional and explicit. Each call records
the provider, purpose, input class, retention setting, and approval. A
project may forbid remote calls entirely; the deterministic path must still be
useful.

The system stores hashes and minimum necessary excerpts in the evaluation
record. It should not retain a full manuscript in a telemetry stream merely
because a finding was emitted. A reader study may require the rendered
document, but that is a declared study object with a retention decision, not an
unnoticed side effect of a linter.

\section{The first release and its boundary}

The first production release is complete when one finished LaTeX manuscript can
travel from an immutable snapshot to a reader packet through the full
humanization path. It includes:

\begin{enumerate}
\item a humanization brief schema and a concept-baseline builder with concept,
      dependency, obligation, and scaffolding fixtures;
\item a parser adapter with protected-span fixtures and source locations;
\item a source-grounded writer interface that accepts only one approved source
      span and its assigned concepts;
\item the integrity, meaning, writing, and reconstruction verification families,
      with stable machine-readable output;
\item a cause-specific repair controller with cycle and oscillation limits;
\item an author decision view, patch assembly, protected comparison, whole-tree
      blackline build, and fail-closed release gate; and
\item a calibration and held-out corpus plus a replay harness for the benchmark
      and held-out manuscripts.
\end{enumerate}

It does not include a general editor, cloud collaboration, enterprise
authentication, autonomous acceptance, authorship classification, or a
population-level style dashboard. It also does not promise a subjective prose
certificate. Those may be sensible later products, but including them now would
make the causal question and the implementation boundary harder to see.

\section{Failure handling}

The system must fail visibly. If parsing loses a macro, it returns an
unparsed-source status and blocks release. If a citation identity cannot be
resolved, it enters a human queue rather than being silently accepted. If a
remote call is unavailable, the local diagnostics continue but the missing
result remains an abstention that blocks any gate that depends on it. If a
protected object changes, the author must explain the change before the reader
packet can be released. If a repair cycle repeats a failure, the controller
stops rather than polishing the same defect indefinitely.

These are engineering controls, not a substitute for scholarship. A passing
fixture proves that the implementation handled that fixture. It does not prove
that the resulting prose is clear or that a reader will agree with the
proposal. The evaluation chapter keeps those judgments separate.

\section{The build path at a glance}
\label{sec:core-acceptance}

The canonical acceptance register appears once in
Appendix~\ref{app:requirements}. The implementation contract in
Appendix~\ref{app:implementation-contract} defines how those requirements
are executed, and \path{schemas/} provides the machine-readable records. The
shorter table below is a reader-facing map of the build path; it deliberately
does not create a second set of release tests.

\begin{table}[H]
\centering
\small
\begin{tabularx}{\textwidth}{@{}p{0.27\textwidth}L p{0.27\textwidth}@{}}
\toprule
Stage & What the team produces & What happens if it fails \\
\midrule
Brief and frozen baseline & The reader, evidence boundary, every substantive
source concept, its explanation obligations, and the dependency order are
explicit before any prose changes. & Stop before writer invocation and ask the
owner to resolve the missing field or ambiguity. \\
Protected core & A LaTeX source snapshot, canonical PDF, source map, and
protected comparison can be replayed by a second operator. & Narrow the release
to the protected review utility. \\
Source-grounded rewriting & One approved passage can be rewritten, rebuilt, and
repaired without losing a concept or widening a claim. & Keep model-dependent
work in abstention and return to the deterministic path. \\
Verification and release & Meaning, integrity, writing, evidence, and
reconstruction checks produce a located decision before a reader packet exists. &
Repair or stop; never use the reader as the ordinary debugger. \\
Reader packet & A named reader receives one clean PDF, its blackline comparison,
the decision question, and a short reconstruction task. & Preserve the source and
report noncompletion. \\
Product evidence & The sponsor receives benchmark, held-out, replay, burden,
safety, rights, and cost evidence. & Do not make an efficacy or scale claim. \\
\bottomrule
\end{tabularx}
\caption{The reader-facing map of the implementation path. The canonical
acceptance register is in Appendix~\ref{app:requirements} (requirements R1--R29);
exit codes and execution boundaries are in the implementation contract annex.
This table provides a high-level narrative view; R1--R29 is the authoritative
release gate.}
\label{tab:build-path-summary}
\end{table}

The live feasibility document starts only after these tests pass. If the tests
pass and the reader still cannot state the purpose or requested action, the
workflow has failed its product purpose even though its software is running.
